% Source: physical PDF pp. 183--198 (printed pp. 160--175), from the Section
% 27.9 heading on physical p. 183 through the final paragraph immediately
% before Problems on physical p. 198.
% Inventory: Eqs. (27.9.1)--(27.9.52), six ordinary footnotes, no figures,
% tables, or typographic dividers.
% Source notes: Eq. (27.9.33) visibly leaves the index A unsummed in its two
% terms containing t'_A and mu. Eq. (27.9.50) visibly prints g^2, rather than
% the g^3 suggested by substituting Eqs. (27.9.47) and (27.9.49) into
% Eq. (27.9.45). Both printed forms are preserved.
% Uncertainties: none.

\subsection[Gauge Theories with Extended Supersymmetry]{Gauge Theories with
Extended Supersymmetry\footnotemark}
\footnotetext{This section lies somewhat out of the book's main line of
development and may be omitted in a first reading.}

Theories with unbroken extended supersymmetry are not considered to be good
candidates for realistic extensions of the standard model, because of the
non-chirality of particle multiplets discussed in Section 25.4. Nevertheless
gauge theories with extended supersymmetry are worth some consideration here
because they have provided paradigms for the use of powerful mathematical
methods to solve dynamical problems.

There are a number of special formalisms that have been proposed to construct
Lagrangians with \(N=2\) extended
supersymmetry~\hyperref[ch27-ref-12]{[12]}, but fortunately we can get by with
the tools already at hand. Any theory with \(N=2\) supersymmetry also has
\(N=1\) supersymmetry, so its Lagrangian must be a special case of the
Lagrangians already considered in this chapter. To construct a Lagrangian with
\(N=2\) supersymmetry for some set of the \(N=2\) supermultiplets of particles
constructed in Sections 25.4 and 25.5, we need only write down the most general
Lagrangian with \(N=1\) supersymmetry whose \(N=1\) supermultiplets contain
physical fields for the particles in the \(N=2\) supermultiplets, and then
impose a discrete \(R\)-symmetry on the Lagrangian: a symmetry that acts
differently on different components of \(N=2\) supermultiplets. The Lagrangian
density will then be invariant under a second supersymmetry, whose
supermultiplets are given by acting on the supermultiplets of ordinary \(N=1\)
supersymmetry with the \(R\)-symmetry.

It will be convenient to choose the discrete \(R\)-transformation so that
\begin{equation}
  Q_1\longrightarrow Q_2\,,
  \qquad
  Q_2\longrightarrow-Q_1\,.
  \tag{27.9.1}\label{eq:27.9.1}
\end{equation}
If the central charge were zero then the supersymmetry algebra would be
invariant under an \(SU(2)\) \(R\)-symmetry group, which has the transformation
\eqref{eq:27.9.1} as one finite element \(\exp(i\pi\tau_2/2)\), but symmetry
under the discrete symmetry is sufficient for our purposes, so we do not need
to assume a zero central charge. In fact, it will turn out that the Lagrangians
we construct by this method will have an \(SU(2)\) \(R\)-symmetry, not just
symmetry under the discrete transformation \eqref{eq:27.9.1}.

Let us first consider the renormalizable theory of the gauge bosons of a
general gauge group, together with the superpartners required by \(N=2\)
extended supersymmetry. We saw in Section 25.4 that in \(N=2\) global
supersymmetry theories a massless gauge boson can only belong to a multiplet
also containing a pair of massless fermions of each helicity
\(\mathord{\pm}1/2\) that transform as a doublet under the \(SU(2)\)
\(R\)-symmetry and a pair of \(SU(2)\)-singlet spinless bosons. Since \(N=2\)
supersymmetry includes \(N=1\) supersymmetry, the renormalizable Lagrangian
for this theory must be a special case of the general renormalizable
Lagrangian density \eqref{eq:27.4.1}. One feature of this special case is that
since the gauge boson belongs to the adjoint representation of the gauge
group, so must also the fermions and scalar field. To furnish \(N=2\)
supermultiplets of fields with the correct particle content, we must have one
\(N=1\) chiral superfield \(\Phi_A\) with component fields
\(\phi_A,\psi_A,\mathcal{F}_A\) (with \(\psi_A\) Majorana and \(\phi_A\) and
\(\mathcal{F}_A\) both complex) for each \(N=1\) gauge multiplet
\(V_A^\mu,\lambda_A,D_A\). We impose a discrete \(R\)-symmetry under the
transformation
\begin{equation}
  \psi_A\longrightarrow\lambda_A\,,
  \qquad
  \lambda_A\longrightarrow-\psi_A\,,
  \tag{27.9.2}\label{eq:27.9.2}
\end{equation}
(with all other fields invariant) because this is the effect of the
transformation \eqref{eq:27.9.1}. Since a non-trivial superpotential would
give \(\psi_A\) interactions or mass terms that are absent for \(\lambda_A\),
the superpotential must vanish. The Lagrangian \eqref{eq:27.4.1} therefore
takes the special form
\begin{equation}
  \begin{aligned}
  \mathcal{L}={}&
  -\sum_A(D_\mu\phi)_A^*(D^\mu\phi)_A
  -\frac12\sum_A\left(\bar\psi_A(\sl{D}\psi)_A\right)
  +\sum_A\mathcal{F}_A^*\mathcal{F}_A
  \\
  &-2\sqrt{2}\operatorname{Re}\sum_{ABC}
    C_{ABC}\left(\lambda_{AL}^{\mathrm T}\epsilon\psi_{CL}\right)\phi_B^*
  +i\sum_{ABC}C_{ABC}\phi_B^*\phi_CD_A
  -\sum_A\xi_AD_A+\frac12\sum_AD_AD_A
  \\
  &-\frac14\sum_Af_{A\mu\nu}f_A^{\mu\nu}
  -\frac12\sum_A\left(\bar\lambda_A(\sl{D}\lambda)_A\right)
  +\frac{g^2\theta}{64\pi^2}\epsilon_{\mu\nu\rho\sigma}
    \sum_Af_A^{\mu\nu}f_A^{\rho\sigma}\,,
  \end{aligned}
  \tag{27.9.3}\label{eq:27.9.3}
\end{equation}
where
\begin{equation}
  (D_\mu\psi)_A
  =\partial_\mu\psi_A+\sum_{BC}C_{ABC}V_{B\mu}\psi_C\,,
  \tag{27.9.4}\label{eq:27.9.4}
\end{equation}
\begin{equation}
  (D_\mu\lambda)_A
  =\partial_\mu\lambda_A+\sum_{BC}C_{ABC}V_{B\mu}\lambda_C\,,
  \tag{27.9.5}\label{eq:27.9.5}
\end{equation}
\begin{equation}
  (D_\mu\phi)_A
  =\partial_\mu\phi_A+\sum_{BC}C_{ABC}V_{B\mu}\phi_C\,,
  \tag{27.9.6}\label{eq:27.9.6}
\end{equation}
and
\begin{equation}
  f_{A\mu\nu}
  =\partial_\mu V_{A\nu}-\partial_\nu V_{A\mu}
  +\sum_{BC}C_{ABC}V_{B\mu}V_{C\nu}\,.
  \tag{27.9.7}\label{eq:27.9.7}
\end{equation}
(Recall that in the adjoint representation
\((t_A)_{BC}=-iC_{ABC}\), where \(C_{ABC}\) is the real structure constant,
defined as usual in this book to include factors of gauge couplings, and
taken in a basis in which it is totally antisymmetric.) The Lagrangian
density \eqref{eq:27.9.3} has \(N=1\) supersymmetry with multiplets
\(\phi_A,\psi_A,\mathcal{F}_A\) and \(V_A^\mu,\lambda_A,D_A\), because it is a
special case of Eq.~\eqref{eq:27.4.1}, and it has an \(SU(2)\) symmetry
relating \(\psi_A\) and \(\lambda_A\), including invariance under the finite
\(SU(2)\) transformation \eqref{eq:27.9.2}, so it also has a \emph{second}
independent \(N=1\) supersymmetry with multiplets
\(\phi_A,\lambda_A,\mathcal{F}_A\) and \(V_A^\mu,-\psi_A,D_A\). It therefore
satisfies the conditions imposed by \(N=2\) supersymmetry.

We can eliminate the auxiliary fields by setting them equal to the values at
which the Lagrangian density \eqref{eq:27.9.3} is stationary:
\begin{equation}
  \mathcal{F}_A=0\,,
  \qquad
  D_A=-i\sum_{BC}C_{ABC}\phi_B^*\phi_C\,.
  \tag{27.9.8}\label{eq:27.9.8}
\end{equation}
(We are now assuming that the Fayet--Iliopoulos constants \(\xi_A\) all
vanish.) Inserting these values back in Eq.~\eqref{eq:27.9.3} gives an
equivalent Lagrangian density
\begin{equation}
  \begin{aligned}
  \mathcal{L}={}&
  -\sum_A(D_\mu\phi)_A^*(D^\mu\phi)_A
  -\frac12\sum_A\left(\bar\psi_A(\sl{D}\psi)_A\right)
  \\
  &+\sqrt{2}\sum_{ABC}C_{ABC}
    \left(\bar\psi_B\frac{1-\gamma_5}{2}\lambda_A\right)\phi_C
  -\sqrt{2}\sum_{ABC}C_{ABC}
    \left(\bar\lambda_A\frac{1+\gamma_5}{2}\psi_C\right)\phi_B^*
  -V(\phi,\phi^*)
  \\
  &-\frac14\sum_Af_{A\mu\nu}f_A^{\mu\nu}
  -\frac12\sum_A\left(\bar\lambda_A(\sl{D}\lambda)_A\right)
  +\frac{g^2\theta}{64\pi^2}\epsilon_{\mu\nu\rho\sigma}
    \sum_Af_A^{\mu\nu}f_A^{\rho\sigma}\,,
  \end{aligned}
  \tag{27.9.9}\label{eq:27.9.9}
\end{equation}
where the potential is
\begin{equation}
  V(\phi,\phi^*)
  =-\frac12\sum_A\left[\sum_{BC}C_{ABC}\phi_B^*\phi_C\right]^2
  =2\sum_A\left[\sum_{BC}C_{ABC}\operatorname{Re}\phi_B
    \operatorname{Im}\phi_C\right]^2\,.
  \tag{27.9.10}\label{eq:27.9.10}
\end{equation}
This potential has a minimum value of zero, which is reached not only for
\(\phi_A=0\) but also for any set of \(\phi\)s for which
\(\sum_{BC}C_{ABC}\phi_B^*\phi_C=0\) for all \(A\) or, in other words, for
which
\begin{equation}
  [t\mathbin{\cdot}\operatorname{Re}\phi,
   t\mathbin{\cdot}\operatorname{Im}\phi]=0\,,
  \qquad
  \text{where}\qquad
  t\mathbin{\cdot}v\equiv\sum_Bt_Bv_B\,.
  \tag{27.9.11}\label{eq:27.9.11}
\end{equation}
That is, the minimum of the potential is reached for those scalar fields for
which all generators \(t\mathbin{\cdot}\operatorname{Re}\phi\) and
\(t\mathbin{\cdot}\operatorname{Im}\phi\) belong to a Cartan subalgebra of
the full gauge algebra, all of whose generators commute with one another.
Though all such values of \(\phi\) give zero potential, and hence unbroken
\(N=2\) supersymmetry, they are not physically equivalent, as shown for
instance by the different masses they give the gauge bosons associated with
the broken gauge symmetries.

One remarkable feature of extended supersymmetry is that the central charges
of the supersymmetry algebra in any state can be calculated in terms of
`charges' in that state to which bosonic fields are
coupled~\hyperref[ch27-ref-13]{[13]}. The easiest way to do this calculation
is to use the transformation properties under ordinary \(N=1\)
supersymmetry of the extended supersymmetry currents \(S_r^\mu(x)\) with
\(r=2,3,\ldots,N\) to calculate the anticommutators
\(\{Q_{1\alpha},S_{r\beta}^\mu(x)\}\). We can then calculate the central
charges from the anticommutators
\begin{equation}
  \{Q_{1\alpha},Q_{r\beta}\}
  =\int d^3x\,\{Q_{1\alpha},S_{r\beta}^0(x)\}\,.
  \tag{27.9.12}\label{eq:27.9.12}
\end{equation}
It turns out that the integrand on the right-hand side is a derivative with
respect to space coordinates, but its integral does not vanish if the states
have fields that do not vanish rapidly as \(\mathbf{x}\longrightarrow\infty\).

To see how this works in detail, let's consider the case of \(N=2\)
supersymmetry with an \(SU(2)\) gauge symmetry and a single \(N=2\) gauge
supermultiplet, with no additional matter supermultiplets. The Lagrangian here
is given by Eq.~\eqref{eq:27.9.3}, with \(A\), \(B\), and \(C\) running over
the values 1, 2, 3, and
\begin{equation}
  C_{ABC}=e\epsilon_{ABC}\,,
  \qquad
  \xi_A=0\,.
  \tag{27.9.13}\label{eq:27.9.13}
\end{equation}
(The coupling constant here is denoted \(e\) because this is the charge with
which the massless gauge field of the unbroken \(U(1)\) gauge symmetry
interacts.) The usual \(N=1\) supersymmetry current (distinguished now with a
subscript 1) is given by Eq.~\eqref{eq:27.4.40} as
\begin{equation}
  \begin{aligned}
  S_1^\mu={}&
  -\frac14\sum_A f_{A\rho\sigma}
    [\gamma^\rho,\gamma^\sigma]\gamma^\mu\lambda_A
  -e\sum_{ABC}\epsilon_{ABC}\gamma_5\gamma^\mu
    \lambda_A\phi_B^*\phi_C
  \\
  &+\frac{1}{\sqrt{2}}\sum_A
    \left[
      (\sl{D}\phi)_A\gamma^\mu\psi_{AR}
      +(\sl{D}\phi^*)_A\gamma^\mu\psi_{AL}
    \right].
  \end{aligned}
  \tag{27.9.14}\label{eq:27.9.14}
\end{equation}
We can calculate the second supersymmetry current by subjecting \(S_1^\mu\)
to the finite \(SU(2)\) \(R\)-symmetry used above, which simply amounts to
making the replacements
\(\psi_A\longrightarrow\lambda_A\),
\(\lambda_A\longrightarrow-\psi_A\). This gives
\begin{equation}
  \begin{aligned}
  S_2^\mu={}&
  \frac14\sum_A f_{A\rho\sigma}
    [\gamma^\rho,\gamma^\sigma]\gamma^\mu\psi_A
  +e\sum_{ABC}\epsilon_{ABC}\gamma_5\gamma^\mu
    \psi_A\phi_B^*\phi_C
  \\
  &+\frac{1}{\sqrt{2}}\sum_A
    \left[
      (\sl{D}\phi)_A\gamma^\mu\lambda_{AR}
      +(\sl{D}\phi^*)_A\gamma^\mu\lambda_{AL}
    \right].
  \end{aligned}
  \tag{27.9.15}\label{eq:27.9.15}
\end{equation}
It will be enough for our purposes (and somewhat easier) to calculate only
the change in the right-handed part of this current under an \(N=1\)
supersymmetry transformation. After setting the auxiliary fields equal to
their equilibrium values
\[
  \mathcal{F}_A=0\,,
  \qquad
  D_A=-ie\sum_{BC}\epsilon_{ABC}\phi_B^*\phi_C\,,
\]
we find
\[
  \begin{aligned}
  \delta S_{2R}^\mu={}&
  \frac{\sqrt{2}}{4}\sum_A f_{A\rho\sigma}
    [\gamma^\rho,\gamma^\sigma]\gamma^\mu
    (\sl{D}\phi)_A\alpha_R
  \\
  &-\sqrt{2}e\sum_{ABC}\epsilon_{ABC}\gamma^\mu
    (\sl{D}\phi)_A\alpha_R\phi_B^*\phi_C
  -\frac{\sqrt{2}}{4}f_{A\rho\sigma}(\sl{D}\phi)_A
    \gamma^\mu[\gamma^\rho,\gamma^\sigma]\alpha_R
  \\
  &-\sqrt{2}e\sum_{ABC}\epsilon_{ABC}
    \phi_B^*\phi_C(\sl{D}\phi)_A\gamma^\mu\alpha_R+\cdots\,,
  \end{aligned}
\]
where the dots denote terms bilinear in fermion fields, which do not concern
us here because we are interested in effects of long range boson fields. We
can combine terms by using the Dirac anticommutation relations and the
identity
\[
  [\gamma^\rho,\gamma^\sigma][\gamma^\mu,\gamma^\nu]
  +[\gamma^\mu,\gamma^\nu][\gamma^\rho,\gamma^\sigma]
  =-8\eta^{\mu\rho}\eta^{\nu\sigma}
  +8\eta^{\sigma\mu}\eta^{\rho\nu}
  +8i\epsilon^{\mu\nu\rho\sigma}\gamma_5
\]
and find
\[
  \begin{aligned}
  \delta S_{2R}^\mu={}&
  -2\sqrt{2}\sum_Af_A^{\mu\nu}(D_\nu\phi)_A\alpha_R
  -i\sqrt{2}\sum_A\epsilon^{\mu\nu\rho\sigma}
    f_{A\rho\sigma}(D_\nu\phi)_A\alpha_R
  \\
  &-2\sqrt{2}e\sum_{ABC}\epsilon_{ABC}\phi_B^*\phi_C
    (D^\mu\phi)_A\alpha_R+\cdots\,.
  \end{aligned}
\]
In order to write this as a derivative, we need to use the Yang--Mills field
equations given by Eqs.~(15.3.6), (15.3.7), and (15.3.9):
\[
  \begin{aligned}
  D_\nu f_A^{\mu\nu}
  =J_A^\mu
  &=e\sum_{BC}\epsilon_{ABC}
    \left((D^\mu\phi)_B^*\phi_C-\phi_B^*(D^\mu\phi)_C\right),
  \\
  \epsilon_{\mu\nu\rho\sigma}(D^\nu f^{\rho\sigma})_A&=0\,.
  \end{aligned}
\]
These allow us to write \(\delta S_{2R}^\mu\) as a total derivative
\begin{equation}
  \delta S_{2R}^\mu=D_\nu X^{\mu\nu}\alpha_R\,,
  \tag{27.9.16}\label{eq:27.9.16}
\end{equation}
where
\begin{equation}
  X^{\mu\nu}
  =-2\sqrt{2}\sum_Af_A^{\mu\nu}\phi_A
  -i\sqrt{2}\sum_A\epsilon^{\mu\nu\rho\sigma}
    f_{A\rho\sigma}\phi_A+\cdots\,,
  \tag{27.9.17}\label{eq:27.9.17}
\end{equation}
with the dots again indicating irrelevant terms involving fermion fields.
Eq.~\eqref{eq:26.1.18} allows us to write Eq.~\eqref{eq:27.9.16} as an
anticommutation relation
\begin{equation}
  \{Q_{R\alpha},S_{R\beta}^\mu\}
  =i\left[\epsilon\left(\frac{1-\gamma_5}{2}\right)\right]_{\alpha\beta}
    D_\nu X^{\mu\nu}\,.
  \tag{27.9.18}\label{eq:27.9.18}
\end{equation}
Since \(X^{\mu\nu}\) is a gauge-invariant quantity, its gauge-covariant
derivative is the same as its ordinary derivative. Also, \(X^{\mu\nu}\) is
antisymmetric, so \(D_\nu X^{0\nu}=\partial_iX^{0i}\). From
Eqs.~\eqref{eq:27.9.12} and \eqref{eq:27.9.18}, we have at last
\begin{equation}
  \{Q_{R\alpha},Q_{R\beta}\}
  =i\left[\epsilon\left(\frac{1-\gamma_5}{2}\right)\right]_{\alpha\beta}
    \int dS_i\,X^{0i}\,,
  \tag{27.9.19}\label{eq:27.9.19}
\end{equation}
with the integral taken over a large closed surface enclosing the system in
question, with surface area differential \(dS\) taken normal to the surface.
Comparing this with Eq.~\eqref{eq:25.2.38} gives the central charge
\begin{equation}
  Z_{12}=-i\int dS_i\,X^{0i}\,.
  \tag{27.9.20}\label{eq:27.9.20}
\end{equation}
If we choose a gauge in which \(\phi_A\) (almost everywhere) has only the
constant non-vanishing component \(\phi_3=v\), then
\begin{equation}
  \sum_Af_A^{0i}\phi_A=-vE^i\,,
  \qquad
  \frac12\sum_A\epsilon^{0i\rho\sigma}f_{A\rho\sigma}\phi_A=vB^i\,,
  \tag{27.9.21}\label{eq:27.9.21}
\end{equation}
where \(\mathbf{E}\) and \(\mathbf{B}\) are the electric and magnetic fields
associated with the unbroken \(U(1)\) subgroup of the \(SU(2)\) gauge group.
Therefore the central charge \eqref{eq:27.9.20} here is
\begin{equation}
  Z_{12}=2\sqrt{2}v[iq-\mathcal{M}]\,,
  \tag{27.9.22}\label{eq:27.9.22}
\end{equation}
where \(q\) and \(\mathcal{M}\) are the electric charge and magnetic monopole
moment, defined by
\begin{equation}
  q=\int dS_i\,E^i\,,
  \qquad
  \mathcal{M}=\int dS_i\,B^i\,.
  \tag{27.9.23}\label{eq:27.9.23}
\end{equation}
As discussed in Section 23.3, this theory, with \(SU(2)\) gauge symmetry
spontaneously broken by the expectation value of an \(SU(2)\) triplet of
scalars, is one in which magnetic monopoles actually do occur.

The application of the results of Section 27.4 to the Lagrangian density
\eqref{eq:27.9.3} shows that, after the spontaneous breaking of the \(SU(2)\)
gauge symmetry, this theory will contain elementary particles of charge
\(\mathord{\pm}e\), zero magnetic monopole moment, and tree-approximation mass
\(M=\sqrt{2}|ev|\). Specifically, for each sign of the charge there is one
such particle of spin 1, two of them with spin \(1/2\), and one with spin 0.
It is a striking consequence of the results obtained here that \emph{the mass
value \(\sqrt{2}|ev|\) is exact, being unaffected by either radiative
corrections or non-perturbative effects, provided that the quantity \(v\) is
defined by Eq.~\eqref{eq:27.9.22} for the central
charge}~\hyperref[ch27-ref-13]{[13]}.

To see this, note that the massive one-particle states for each sign of the
charge are `short' \(N=2\) supermultiplets, which, as shown at the end of
Section 25.5, have masses that saturate the lower bound \eqref{eq:25.5.24}:
\begin{equation}
  M=|Z_{12}|/2\,.
  \tag{27.9.24}\label{eq:27.9.24}
\end{equation}
Even if we did not trust the tree approximation to give us the precise value
of the particle masses, we would not expect corrections to this approximation
to turn short multiplets into the full multiplets, with many more states, that
could have larger masses, so we can be confident that Eq.~\eqref{eq:27.9.24}
is exactly valid. For particles with electric charge \(q=\mathord{\pm}e\) and
zero magnetic monopole moment, Eq.~\eqref{eq:27.9.20} gives
\(Z_{12}=\mathord{\pm}2\sqrt{2}ive\), so Eq.~\eqref{eq:27.9.24} tells us that
their masses are
\begin{equation}
  M=\sqrt{2}|ev|\,.
  \tag{27.9.25}\label{eq:27.9.25}
\end{equation}
This is the result found in the tree approximation, but now we see that it is
exact.

The semi-classical calculations described in Section 23.3 show that
electrically neutral magnetic monopoles in this theory have monopole
strengths\footnote{Note that the magnetic moment \(\mathcal{M}\) defined by
Eq.~\eqref{eq:27.9.23} is related to the magnetic moment \(g\) defined in
Section 23.3 by \(\mathcal{M}=4\pi g\).}
\begin{equation}
  \mathcal{M}=\frac{4\pi\nu}{e}\,,
  \tag{27.9.26}\label{eq:27.9.26}
\end{equation}
with \(\nu\) the winding number, a positive or negative integer. The formula
\eqref{eq:27.9.22} for the central charge together with the inequality
\eqref{eq:25.5.24} thus give a lower bound on the monopole masses
\begin{equation}
  M\geq\frac{4\pi\sqrt{2}\,|\nu v|}{|e|}\,.
  \tag{27.9.27}\label{eq:27.9.27}
\end{equation}
Interestingly, this is the same as the Bogomol'nyi lower
bound~\hyperref[ch27-ref-14]{[14]} on monopole energies derived in Section
23.3.\footnote{The canonically normalized field with a non-vanishing vacuum
expectation value is \(\sqrt{2}\operatorname{Re}\phi_3\) (for real \(v\)), so
the quantity \(\langle\phi\rangle\) appearing in the Bogomol'nyi inequality
(23.3.19) is \(\sqrt{2}v\).} In fact, the monopole solution for \(\nu=1\)
that was described in Section 23.3 saturates this bound. More generally, the
`dyons' of this theory~\hyperref[ch27-ref-15]{[15]}, particles with both
charge and magnetic moment, have masses given
by~\hyperref[ch27-ref-16]{[16]}
\begin{equation}
  M=2|v|\sqrt{q^2+\mathcal{M}^2}\,,
  \tag{27.9.28}\label{eq:27.9.28}
\end{equation}
which again is the minimum value allowed by Eqs.~\eqref{eq:25.5.24} and
\eqref{eq:27.9.20}. Indeed, all of the known particles in this theory have
masses given in the semi-classical limit by
Eq.~\eqref{eq:27.9.28}~\hyperref[ch27-ref-17]{[17]}.

Returning now to general \(N=2\) gauge theories, we may also include
additional `matter' fields in the Lagrangian. For simplicity, we will restrict
ourselves to `short' massive hypermultiplets (with central charge
\(\mathcal{Z}\) saturating the inequality \eqref{eq:25.5.24}), each consisting
of a single fermion of spin \(1/2\) and an \(SU(2)\) doublet of spin 0
particles, together with distinct antiparticles. This is the same spin
content as is given under \(N=1\) supersymmetry by pairs of left-chiral
scalar superfields \(\Phi'_n\) and \(\Phi''_n\), together with their
right-chiral adjoints, with the complex scalar field components \(\phi'_n\)
and \(\phi''_n\) and their adjoints forming pairs of \(SU(2)\) doublets, and
the spinor fields all \(SU(2)\) singlets. (We are using primes and double
primes to distinguish these superfields and their components from \(\Phi_A\)
and its components.) If some of these hypermultiplets \(\Phi'_n\) and
\(\Phi''_n\) are non-neutral under the gauge group, then a superpotential is
allowed, of the form:\footnote{There still cannot be any terms in the
superpotential that are of second or higher order in the \(\Phi_A\) for the
same reason as before: such terms would lead to scalar couplings or masses
for the \(\psi_A\), with no corresponding couplings or masses for their
\(SU(2)\) partners \(\lambda_A\). Also, there cannot be any term that is
trilinear in the \(\Phi'_n\) and/or \(\Phi''_n\), because then the term in
Eq.~\eqref{eq:27.4.1} involving the product of fermion bilinears with second
derivatives of the superpotential would lead to a coupling of the
\(SU(2)\)-singlet fermions to \(SU(2)\) doublet fields \(\phi'_n\) or
\(\phi''_n\). Thus the only trilinear interaction terms in the
superpotential must involve one factor of \(\Phi_A\) and two factors of
\(\Phi'_n\) and/or \(\Phi''_n\). There cannot be any trilinear terms
involving a \(\Phi_A\) and two \(\Phi'_n\)s or two \(\Phi''_n\)s, because
that would give the \(SU(2)\) singlet auxiliary fields \(\mathcal{F}_A\) an
interaction with \(SU(2)\)-triplet products \(\phi'_n\phi'_m\) or
\(\phi''_n\phi''_m\), and there cannot be any bilinear terms involving two
\(\Phi'_n\)s or two \(\Phi''_n\)s, because that would yield \(SU(2)\)-triplet
mass terms \((\psi_{nL}^{\prime\mathrm T}\epsilon\psi'_{mL})\) or
\((\psi_{nL}^{\prime\prime\mathrm T}\epsilon\psi''_{mL})\). The only
remaining allowed bilinear or trilinear terms are of the form
\eqref{eq:27.9.29}.}
\begin{equation}
  f(\Phi,\Phi',\Phi'')
  =\frac12\sum_{Anm}(s_A)_{nm}\Phi'_n\Phi''_m\Phi_A
  +\frac12\sum_{nm}\mu_{nm}\Phi'_n\Phi''_m\,.
  \tag{27.9.29}\label{eq:27.9.29}
\end{equation}
To the Lagrangian density \eqref{eq:27.9.3} we then must add a Lagrangian
density for these hypermultiplets, given by the first eight terms on the
right-hand side of Eq.~\eqref{eq:27.4.1}, and find a total Lagrangian density
\begin{equation}
  \begin{aligned}
  \mathcal{L}={}&
  -\sum_n(D_\mu\phi')_n^*(D^\mu\phi')_n
  -\sum_n(D_\mu\phi'')_n^*(D^\mu\phi'')_n
  -\sum_A(D_\mu\phi)_A^*(D^\mu\phi)_A
  \\
  &-\frac12\sum_n\left(\bar\psi'_n(\sl{D}\psi')_n\right)
  -\frac12\sum_n\left(\bar\psi''_n(\sl{D}\psi'')_n\right)
  -\frac12\sum_A\left(\bar\psi_A(\sl{D}\psi)_A\right)
  -\frac12\sum_A\left(\bar\lambda_A(\sl{D}\lambda)_A\right)
  \\
  &+\sum_n\mathcal{F}_n^{\prime *}\mathcal{F}'_n
  +\sum_n\mathcal{F}_n^{\prime\prime *}\mathcal{F}''_n
  +\sum_A\mathcal{F}_A^*\mathcal{F}_A
  \\
  &-\operatorname{Re}\sum_{Anm}(s_A)_{nm}\phi_A
    \left(\psi_{nL}^{\prime\mathrm T}\epsilon\psi''_{mL}\right)
  -2\sqrt{2}\operatorname{Re}\sum_{ABC}C_{ABC}
    \left(\lambda_{AL}^{\mathrm T}\epsilon\psi_{CL}\right)\phi_B^*
  \\
  &-\operatorname{Re}\sum_{Anm}(s_A)_{nm}\phi'_n
    \left(\psi_{mL}^{\prime\prime\mathrm T}\epsilon\psi_{AL}\right)
  -\operatorname{Re}\sum_{Anm}(s_A)_{nm}\phi''_m
    \left(\psi_{nL}^{\prime\mathrm T}\epsilon\psi_{AL}\right)
  \\
  &+2\sqrt{2}\operatorname{Im}\sum_{Anm}(t'_A)_{mn}
    \left(\psi_{nL}^{\prime\mathrm T}\epsilon\lambda_{AL}\right)\phi_m^{\prime *}
  +2\sqrt{2}\operatorname{Im}\sum_{Anm}(t''_A)_{mn}
    \left(\psi_{nL}^{\prime\prime\mathrm T}\epsilon\lambda_{AL}\right)
    \phi_m^{\prime\prime *}
  \\
  &+\operatorname{Re}\sum_{Anm}(s_A)_{nm}\phi_A\phi'_n\mathcal{F}''_m
  +\operatorname{Re}\sum_{Anm}(s_A)_{nm}\phi_A\phi''_m\mathcal{F}'_n
  +\operatorname{Re}\sum_{Anm}(s_A)_{nm}\phi'_n\phi''_m\mathcal{F}_A
  \\
  &+\operatorname{Re}\sum_{nm}\mu_{nm}\phi'_n\mathcal{F}''_m
  +\operatorname{Re}\sum_{nm}\mu_{nm}\phi''_m\mathcal{F}'_n
  -\operatorname{Re}\sum_{nm}\mu_{nm}
    \left(\psi_{nL}^{\prime\mathrm T}\epsilon\psi''_{mL}\right)
  \\
  &-\sum_{Anm}(t'_A)_{nm}\phi_n^{\prime *}\phi'_mD_A
  -\sum_{Anm}(t''_A)_{nm}\phi_n^{\prime\prime *}\phi''_mD_A
  +i\sum_{ABC}C_{ABC}\phi_B^*\phi_CD_A
  \\
  &-\sum_A\xi_AD_A+\frac12\sum_AD_AD_A
  -\frac14\sum_Af_{A\mu\nu}f_A^{\mu\nu}
  +\frac{g^2\theta}{64\pi^2}\epsilon_{\mu\nu\rho\sigma}
    \sum_Af_A^{\mu\nu}f_A^{\rho\sigma}\,,
  \end{aligned}
  \tag{27.9.30}\label{eq:27.9.30}
\end{equation}
where the \((t'_A)_{nm}\) and \((t''_A)_{nm}\) are the matrices (including
coupling constant factors) representing the gauge group on the left-chiral
scalar superfields \(\Phi'_n\) and \(\Phi''_n\), respectively. The Yukawa
couplings between fermions and scalars have a discrete \(R\)-symmetry under
the transformation
\begin{equation}
  \lambda_{AL}\longrightarrow-\psi_{AL}\,,
  \quad
  \psi_{AL}\longrightarrow+\lambda_{AL}\,,
  \quad
  \phi''_n\longrightarrow-\phi_n^{\prime *}\,,
  \quad
  \phi'_n\longrightarrow\phi_n^{\prime\prime *}\,,
  \tag{27.9.31}\label{eq:27.9.31}
\end{equation}
provided that
\begin{equation}
  s_A=-2\sqrt{2}i\,t_A^{\prime\mathrm T}
  =+2\sqrt{2}i\,t''_A\,.
  \tag{27.9.32}\label{eq:27.9.32}
\end{equation}
(Note in particular that Eq.~\eqref{eq:27.9.32} requires the representations
of the gauge group furnished by \(\Phi'_n\) and \(\Phi''_n\) to be complex
conjugates.) This is also a symmetry of all the other terms in the Lagrangian
density \eqref{eq:27.9.30}, except for those involving the auxiliary fields.

It is not possible to extend the symmetry under the transformation
\eqref{eq:27.9.31} to the auxiliary fields, but the symmetry appears after
the auxiliary fields are eliminated.\footnote{After elimination of the
auxiliary fields, the resulting action is invariant under the original
\(N=2\) supersymmetry transformation only `on-shell'---that is, only up to
terms that vanish when the fields satisfy the interacting field equations.
This doesn't hurt, because there still are two conserved supersymmetry
currents whose integrated time-components satisfy the \(N=2\) supersymmetry
anticommutation relations when the fields of which they are composed are
required to satisfy the field equations of the Heisenberg picture. `Off-shell'
formulations of \(N=2\) supersymmetry exist, but are subject to various
complications~\hyperref[ch27-ref-18]{[18]}.} After setting
\(D_A,\mathcal{F}'_n\), and \(\mathcal{F}''_n\) equal to values at which the
Lagrangian density is stationary, and combining \(D\)- and
\(\mathcal{F}\)-terms, the Lagrangian density (with \(s_A\) and \(t''_A\)
given by Eq.~\eqref{eq:27.9.32}, and \(\xi_A\) taken to vanish) takes the form
\begin{equation}
  \begin{aligned}
  \mathcal{L}={}&
  -\sum_n(D_\mu\phi')_n^*(D^\mu\phi')_n
  -\sum_n(D_\mu\phi'')_n^*(D^\mu\phi'')_n
  -\sum_A(D_\mu\phi)_A^*(D^\mu\phi)_A
  \\
  &-\frac12\sum_n\left(\bar\psi'_n(\sl{D}\psi')_n\right)
  -\frac12\sum_n\left(\bar\psi''_n(\sl{D}\psi'')_n\right)
  -\frac12\sum_A\left(\bar\psi_A(\sl{D}\psi)_A\right)
  -\frac12\sum_A\left(\bar\lambda_A(\sl{D}\lambda)_A\right)
  \\
  &-2\sqrt{2}\operatorname{Im}\sum_{Anm}(t'_A)_{mn}\phi_A
    \left(\psi_{nL}^{\prime\mathrm T}\epsilon\psi''_{mL}\right)
  -2\sqrt{2}\operatorname{Re}\sum_{ABC}C_{ABC}
    \left(\lambda_{AL}^{\mathrm T}\epsilon\psi_{CL}\right)\phi_B^*
  \\
  &-2\sqrt{2}\operatorname{Im}\sum_{Anm}(t'_A)_{mn}\phi'_n
    \left(\psi_{mL}^{\prime\prime\mathrm T}\epsilon\psi_{AL}\right)
  -2\sqrt{2}\operatorname{Im}\sum_{Anm}(t'_A)_{mn}\phi''_m
    \left(\psi_{nL}^{\prime\mathrm T}\epsilon\psi_{AL}\right)
  \\
  &+2\sqrt{2}\operatorname{Im}\sum_{Anm}(t'_A)_{mn}
    \left(\psi_{nL}^{\prime\mathrm T}\epsilon\lambda_{AL}\right)\phi_m^{\prime *}
  -2\sqrt{2}\operatorname{Im}\sum_{Anm}(t'_A)_{nm}
    \left(\psi_{nL}^{\prime\prime\mathrm T}\epsilon\lambda_{AL}\right)
    \phi_m^{\prime\prime *}
  \\
  &-\frac14\sum_Af_{A\mu\nu}f_A^{\mu\nu}
  +\frac{g^2\theta}{64\pi^2}\epsilon_{\mu\nu\rho\sigma}
    \sum_Af_A^{\mu\nu}f_A^{\rho\sigma}
  \\
  &-\sum_{ABnm}\{t'_A,t'_B\}_{mn}\phi_A\phi_B^*
    \left(\phi_n^{\prime *}\phi'_m+\phi_n^{\prime\prime *}\phi''_m\right)
  \\
  &-\frac12\sum_A\left[
    \sum_{nm}(t'_A)_{nm}
    \left(\phi_n^{\prime *}\phi'_m-\phi''_n\phi_m^{\prime\prime *}\right)
  \right]^2
  \\
  &+\frac12\sum_{ABCDE}C_{ABC}C_{ADE}
    \phi_B^*\phi_C\phi_D^*\phi_E
  -2\sum_A\left|\sum_{nm}(t'_A)_{nm}\phi'_n\phi''_m\right|^2
  \\
  &-4\operatorname{Re}\sum_{nm}(t'_A\mu)_{nm}\phi_n^{\prime *}\phi'_m
  -4\operatorname{Re}\sum_{nm}(\mu t'_A)_{nm}
    \phi''_n\phi_m^{\prime\prime *}
  \\
  &-2\sum_{nm}(\mu^\dagger\mu)_{nm}\phi_n^{\prime *}\phi'_m
  -2\sum_{nm}(\mu\mu^\dagger)_{nm}\phi''_n\phi_m^{\prime\prime *}\,.
  \end{aligned}
  \tag{27.9.33}\label{eq:27.9.33}
\end{equation}
The last five lines on the right-hand side come from the terms in
Eq.~\eqref{eq:27.9.30} involving auxiliary fields, and now these too are
invariant under the discrete transformation \eqref{eq:27.9.31}, provided
that
\begin{equation}
  [t'_A,\mu]=[\mu^\dagger,\mu]=0\,.
  \tag{27.9.34}\label{eq:27.9.34}
\end{equation}

We can now go a step further and consider the case of \(N=4\) extended global
supersymmetry. (As remarked in Section 25.4, \(N=3\) supersymmetry is the same
as \(N=4\) supersymmetry.) The only massless multiplets of \(N=4\)
supersymmetry that do not contain gravitons or gravitinos consist of a single
particle of helicity 1, an \(SU(4)\) quartet of particles of helicity \(1/2\),
and an \(SU(4)\) sextet of particles of helicity 0, together with their
CPT-conjugates of opposite helicity. There is one such supermultiplet for
each generator \(t_A\) of the gauge group. These particles can be grouped
into supermultiplets of \(N=2\) supersymmetry: for each \(t_A\) there is one
gauge supermultiplet consisting of one particle of helicity 1, two particles
of helicity \(\mathord{\pm}1/2\), and one particle of helicity 0, together
with their CPT-conjugates of opposite helicity, plus two hypermultiplets,
each consisting of one particle of each helicity \(\mathord{\pm}1/2\) and
two particles of helicity 0. The \(N=2\) gauge superfield consists of an
\(N=1\) gauge superfield \(V_A\) and a left-chiral scalar superfield
\(\Phi_A\) and its complex conjugate, while the two \(N=2\) hypermultiplets
consist of two additional left-chiral scalar superfields \(\Phi'_A\) and
\(\Phi''_A\) and their complex conjugates.

Since \(N=4\) supersymmetry includes \(N=2\) supersymmetry, the
Lagrangian density after elimination of the auxiliary fields of \(N=1\)
supersymmetry~\hyperref[ch27-ref-18]{[18]} must be a special case of
Eq.~\eqref{eq:27.9.33}, but with the labels \(n\), \(m\), etc. running over
the indices \(A\), \(B\), \(C\), etc. of the adjoint representation. Also,
the coefficient \(\mu_{nm}\) in the superpotential \eqref{eq:27.9.29} must
vanish here, because otherwise Eq.~\eqref{eq:27.9.33} would contain terms
quadratic in the fermion fields \(\psi'_A\) and \(\psi''_A\), with no
counterpart for their \(N=4\) superpartners \(\lambda_A\) and \(\psi_A\).
Also setting \((t'_A)_{BC}\) equal to the generators \(-iC_{ABC}\) in the
adjoint representation, we find that the Lagrangian density must take the
form
\begin{equation}
  \begin{aligned}
  \mathcal{L}={}&
  -\sum_A(D_\mu\phi')_A^*(D^\mu\phi')_A
  -\sum_A(D_\mu\phi'')_A^*(D^\mu\phi'')_A
  -\sum_A(D_\mu\phi)_A^*(D^\mu\phi)_A
  \\
  &-\frac12\sum_A\left(\bar\psi'_A(\sl{D}\psi')_A\right)
  -\frac12\sum_A\left(\bar\psi''_A(\sl{D}\psi'')_A\right)
  -\frac12\sum_A\left(\bar\psi_A(\sl{D}\psi)_A\right)
  -\frac12\sum_A\left(\bar\lambda_A(\sl{D}\lambda)_A\right)
  \\
  &-2\sqrt{2}\operatorname{Re}\sum_{ABC}C_{ABC}\phi_A
    \left(\psi_{BL}^{\prime\mathrm T}\epsilon\psi''_{CL}\right)
  -2\sqrt{2}\operatorname{Re}\sum_{ABC}C_{ABC}
    \left(\lambda_{AL}^{\mathrm T}\epsilon\psi_{CL}\right)\phi_B^*
  \\
  &-2\sqrt{2}\operatorname{Re}\sum_{ABC}C_{ABC}\phi'_B
    \left(\psi_{CL}^{\prime\prime\mathrm T}\epsilon\psi_{AL}\right)
  -2\sqrt{2}\operatorname{Re}\sum_{ABC}C_{ABC}\phi''_C
    \left(\psi_{BL}^{\prime\mathrm T}\epsilon\psi_{AL}\right)
  \\
  &+2\sqrt{2}\operatorname{Re}\sum_{ABC}C_{ABC}
    \left(\psi_{BL}^{\prime\mathrm T}\epsilon\lambda_{AL}\right)\phi_C^{\prime *}
  +2\sqrt{2}\operatorname{Re}\sum_{ABC}C_{ABC}
    \left(\psi_{BL}^{\prime\prime\mathrm T}\epsilon\lambda_{AL}\right)
    \phi_C^{\prime\prime *}
  \\
  &-\frac14\sum_Af_{A\mu\nu}f_A^{\mu\nu}
  +\frac{g^2\theta}{64\pi^2}\epsilon_{\mu\nu\rho\sigma}
    \sum_Af_A^{\mu\nu}f_A^{\rho\sigma}-V\,,
  \end{aligned}
  \tag{27.9.35}\label{eq:27.9.35}
\end{equation}
where the potential is
\begin{equation}
  \begin{aligned}
  V={}&
  \sum_{ABCDE}C_{ADE}C_{BCE}
    \left(\phi_A\phi_B^*+\phi_B\phi_A^*\right)
    \left(\phi'_C\phi_D^{\prime *}+\phi_C^{\prime\prime *}\phi''_D\right)
  \\
  &+\frac12\sum_A\left|
    \sum_{BC}C_{ABC}
    \left(\phi_B^{\prime *}\phi'_C-\phi''_B\phi_C^{\prime\prime *}\right)
  \right|^2
  \\
  &-\frac12\sum_{ABCDE}C_{ABC}C_{ADE}
    \phi_B^*\phi_C\phi_D^*\phi_E
  +2\sum_A\left|\sum_{BC}C_{ABC}\phi'_B\phi''_C\right|^2\,.
  \end{aligned}
  \tag{27.9.36}\label{eq:27.9.36}
\end{equation}
With no further constraints needed, this Lagrangian has an \(SU(4)\)
\(R\)-symmetry, which implies that it is invariant under \(N=4\)
supersymmetry. To see this, we need to use the Jacobi identity to write the
cross-term in the second line of the right-hand side of Eq.~\eqref{eq:27.9.36}
in the form
\[
  \begin{aligned}
  \sum_{ABCDE}C_{ABC}C_{ADE}
    \phi_B^{\prime *}\phi'_C\phi_D^{\prime\prime *}\phi''_E
  ={}&
  -\sum_{ABCDE}C_{ABC}C_{ADE}
    \phi_B^{\prime *}\phi'_D\phi_E^{\prime\prime *}\phi''_C
  \\
  &-\sum_{ABCDE}C_{ABC}C_{ADE}
    \phi_B^{\prime *}\phi'_E\phi_C^{\prime\prime *}\phi''_D\,,
  \end{aligned}
\]
which allows us to write the potential \eqref{eq:27.9.36} in a form symmetric
among the scalars and their adjoints
\begin{equation}
  \begin{aligned}
  V={}&
  \sum_A\left|\sum_{BC}C_{ABC}\phi_B^*\phi'_C\right|^2
  +\sum_A\left|\sum_{BC}C_{ABC}\phi_B^*\phi_C^{\prime\prime *}\right|^2
  +\sum_A\left|\sum_{BC}C_{ABC}\phi_B\phi'_C\right|^2
  \\
  &+\sum_A\left|\sum_{BC}C_{ABC}\phi_B\phi_C^{\prime\prime *}\right|^2
  +\sum_A\left|\sum_{BC}C_{ABC}\phi_B^{\prime *}\phi''_C\right|^2
  \\
  &+\sum_A\left|\sum_{BC}C_{ABC}\phi'_B\phi''_C\right|^2
  +\frac12\sum_A\left|\sum_{BC}C_{ABC}\phi'_B\phi_C^{\prime *}\right|^2
  \\
  &+\frac12\sum_A\left|\sum_{BC}C_{ABC}\phi''_B
    \phi_C^{\prime\prime *}\right|^2
  +\frac12\sum_A\left|\sum_{BC}C_{ABC}\phi_B\phi_C^*\right|^2\,.
  \end{aligned}
  \tag{27.9.37}\label{eq:27.9.37}
\end{equation}
Now to make the \(SU(4)\) symmetry apparent, we introduce an \(SU(4)\)
notation for the fields. We assemble the left-handed fermion fields into an
\(SU(4)\) vector:
\begin{equation}
  \psi_{1AL}\equiv\psi_{AL}\,,
  \qquad
  \psi_{2AL}\equiv\lambda_{AL}\,,
  \qquad
  \psi_{3AL}\equiv\psi'_{AL}\,,
  \qquad
  \psi_{4AL}\equiv\psi''_{AL}\,.
  \tag{27.9.38}\label{eq:27.9.38}
\end{equation}
In order for the fermion kinematic terms in the Lagrangian density to be
\(SU(4)\)-invariant, we then must assemble the right-handed fermion fields
into a contragredient vector:
\begin{equation}
  \psi^1_{AR}\equiv\psi_{AR}\,,
  \qquad
  \psi^2_{AR}\equiv\lambda_{AR}\,,
  \qquad
  \psi^3_{AR}\equiv\psi'_{AR}\,,
  \qquad
  \psi^4_{AR}\equiv\psi''_{AR}\,.
  \tag{27.9.39}\label{eq:27.9.39}
\end{equation}
The Majorana condition on the fermion fields then takes the
\(SU(4)\)-invariant form
\begin{equation}
  (\psi_{iAL})^*=-\beta\epsilon\psi^i_{AR}\,,
  \tag{27.9.40}\label{eq:27.9.40}
\end{equation}
with indices \(i\), \(j\), etc. running over the values 1, 2, 3, 4. In order
for the Yukawa couplings between the fermion and scalar fields to be
\(SU(4)\)-invariant, we must give the scalars the transformation properties
of an antisymmetric \(SU(4)\) tensor
\begin{equation}
  \begin{gathered}
  \phi_A^{12}\equiv\phi_A^*\,,
  \qquad
  \phi_A^{13}\equiv\phi''_A\,,
  \qquad
  \phi_A^{14}\equiv-\phi'_A\,,
  \\
  \phi_A^{23}\equiv-\phi_A^{\prime *}\,,
  \qquad
  \phi_A^{24}\equiv-\phi_A^{\prime\prime *}\,,
  \qquad
  \phi_A^{34}\equiv\phi_A\,,
  \end{gathered}
  \tag{27.9.41}\label{eq:27.9.41}
\end{equation}
which also obeys an \(SU(4)\)-invariant reality condition
\begin{equation}
  (\phi_A^{ij})^*
  =\frac12\sum_{kl}\epsilon_{ijkl}\phi_A^{kl}\,.
  \tag{27.9.42}\label{eq:27.9.42}
\end{equation}
The whole Lagrangian density \eqref{eq:27.9.35} can then be written in the
manifestly \(SU(4)\)-invariant form
\begin{equation}
  \begin{aligned}
  \mathcal{L}={}&
  -\frac12\sum_{Aij}(D_\mu\phi^{ij})_A(D^\mu\phi^{ij})_A^*
  \\
  &-\frac12\sum_{Ai}
    \left(\psi_{iAL}^{\mathrm T}\epsilon(\sl{D}\psi_R^i)_A\right)
  +\frac12\sum_{Ai}
    \left(\psi_{AR}^{i\mathrm T}\epsilon(\sl{D}\psi_{iL})_A\right)
  \\
  &-\sqrt{2}\operatorname{Re}\sum_{ABCij}C_{ABC}\phi_A^{ij}
    \left(\psi_{iBL}^{\mathrm T}\epsilon\psi_{jCL}\right)-V
  \\
  &-\frac14\sum_Af_{A\mu\nu}f_A^{\mu\nu}
  +\frac{g^2\theta}{64\pi^2}\epsilon_{\mu\nu\rho\sigma}
    \sum_Af_A^{\mu\nu}f_A^{\rho\sigma}\,,
  \end{aligned}
  \tag{27.9.43}\label{eq:27.9.43}
\end{equation}
where the potential is
\begin{equation}
  V=\frac18\sum_{Aijkl}\left|
    \sum_{BC}C_{ABC}\phi_B^{ij}\phi_C^{kl}
  \right|^2\,.
  \tag{27.9.44}\label{eq:27.9.44}
\end{equation}
The potential has a minimum value zero, so that supersymmetry is not broken
in this theory. The minimum is reached when the generators
\(\sum_A t_A\phi_A^{ij}\) all commute with one another.

For zero theta angle, the gauge theories with a simple gauge group and either
\(N=2\) or \(N=4\) supersymmetry have just a single coupling constant, the
gauge coupling constant \(g\). Since these theories have \(N=1\)
supersymmetry, they share the property discussed in Section 27.6, that the
only infinity in higher orders of perturbation theory is in a one-loop
correction to this coupling.\footnote{The trilinear term in the
superpotential \eqref{eq:27.9.29} is proportional to the gauge coupling, and
is therefore renormalized, despite the no-renormalization theorem of Section
27.6. This is because here we are renormalizing the left-chiral scalar
superfields \(\Phi_A,\Phi'_n,\) and \(\Phi''_n\) as well as the gauge
superfield \(V_A\) so as to keep them canonically normalized. The bilinear
term in Eq.~\eqref{eq:27.9.29} is renormalized for the same reason.} The
function \(\beta(g)\) in the renormalization group equation
\(\mu\,dg/d\mu=\beta(g)\) is then given to all orders of perturbation theory
by the one-loop formula (18.7.2), with a suitable correction for the presence
of scalar fields:
\begin{equation}
  \beta(g)
  =-\frac{g^3}{4\pi^2}
    \left(\frac{11}{12}C_1-\frac16C_2^f-\frac1{12}C_2^s\right),
  \tag{27.9.45}\label{eq:27.9.45}
\end{equation}
where
\begin{equation}
  \begin{aligned}
  \sum_{AB}C_{ABC}C_{ABD}
  &=g^2C_1\delta_{CD}\,,
  \\
  \left[\operatorname{Tr}(t_Ct_D)\right]_{\text{Majorana fermions}}
  &=g^2C_2^f\delta_{CD}\,,
  \\
  \left[\operatorname{Tr}(t_Ct_D)\right]_{\text{complex scalars}}
  &=g^2C_2^s\delta_{CD}\,.
  \end{aligned}
  \tag{27.9.46}\label{eq:27.9.46}
\end{equation}
In general theories with \(N=2\) supersymmetry we have two Majorana fermions
\(\lambda_A\) and \(\psi_A\) in the adjoint representation and \(H\) pairs of
Majorana fermions \(\psi'_n\) and \(\psi''_n\) whose left- and right-handed
parts are in representations with generators either \(t'_A\) or
\(-t_A^{\prime\mathrm T}\), so
\begin{equation}
  C_2^f=2C_1+2HC'_2\,,
  \tag{27.9.47}\label{eq:27.9.47}
\end{equation}
where \(C'_2\) is defined by
\begin{equation}
  \operatorname{Tr}t'_Ct'_D=g^2C'_2\delta_{CD}\,.
  \tag{27.9.48}\label{eq:27.9.48}
\end{equation}
Also, we have one complex scalar \(\phi_A\) in the adjoint representation
and \(H\) pairs of complex scalars \(\phi'_n\) and \(\phi''_n\) in
representations with generators \(t'_A\) or \(-t_A^{\prime\mathrm T}\), so
\begin{equation}
  C_2^s=C_1+2HC'_2\,.
  \tag{27.9.49}\label{eq:27.9.49}
\end{equation}
The beta function \eqref{eq:27.9.45} is therefore
\begin{equation}
  \beta(g)=-\frac{g^2}{8\pi^2}(C_1-HC'_2)\,.
  \tag{27.9.50}\label{eq:27.9.50}
\end{equation}
The case of \(N=4\) supersymmetry is just the special case with \(H=1\)
pairs of \(N=2\) hypermultiplets in the adjoint representation, with
\(C'_2=C_1\), so in this case the beta function vanishes. This is therefore
a finite theory, with no renormalizations at
all~\hyperref[ch27-ref-19]{[19]}.

Gauge theories with \(N=4\) supersymmetry have another remarkable property,
known as \emph{duality}. This was first conjectured by Montonen and
Olive~\hyperref[ch27-ref-17]{[17]} for purely bosonic theories in which a
simple gauge group is spontaneously broken to a \(U(1)\) electromagnetic
gauge group. They noticed that semi-classical calculations (of the sort
described in Section 23.3) give the mass of particles with charge \(q=ne\)
and magnetic monopole moment \(\mathcal{M}=4\pi m/e\) (with \(n\) and \(m\)
integers of any sign) as
\begin{equation}
  M=\sqrt{2}\left|v\left(ne+\frac{4\pi im}{e}\right)\right|,
  \tag{27.9.51}\label{eq:27.9.51}
\end{equation}
which is invariant under the transformations
\begin{equation}
  m\longrightarrow n\,,
  \qquad
  n\longrightarrow-m\,,
  \qquad
  e\longleftrightarrow4\pi/e\,.
  \tag{27.9.52}\label{eq:27.9.52}
\end{equation}
On this basis, they suggested that the theory with a weak gauge coupling
\(e\) is fully equivalent to one with a strong gauge coupling \(4\pi/e\).
Neither the purely bosonic theory nor the simplest versions of the \(N=1\)
and \(N=2\) extended supersymmetry theories really have this
property;~\hyperref[ch27-ref-20]{[20]} for one thing, the massive charged
elementary vector bosons of the broken gauge symmetries have spin 1, while
all of the monopoles and dyons have spins \(1/2\) or 0. (We will see in
Section 29.5 that \(N=2\) theories do have a duality property of a more
subtle sort.) But for \(N=4\) supersymmetry the monopole states form
multiplets with one particle of spin 1, four particles of spin \(1/2\), and
two particles of spin 0, just like the elementary
particles~\hyperref[ch27-ref-20]{[20]}. Evidence has
accumulated~\hyperref[ch27-ref-21]{[21]} that \(N=4\) supersymmetric gauge
theories are indeed invariant under the interchange of electric and magnetic
quantum numbers and of \(e\) with \(4\pi/e\). The equivalence of theories
with large and small coupling constants has become an increasingly important
theme in string theory, but this is beyond the scope of this book.
